% crank_nicolson.tex
%
% Crank-Nicolson, the theta-scheme, and Rannacher
% time-stepping.
%
% Self-contained writeup. Compiles standalone with pdflatex.

\documentclass[11pt,a4paper]{article}

\usepackage[T1]{fontenc}
\usepackage[utf8]{inputenc}
\usepackage{amsmath, amssymb, amsthm}
\usepackage{geometry}
\usepackage{hyperref}
\usepackage{enumitem}

\geometry{margin=1in}

\newtheorem{proposition}{Proposition}[section]
\newtheorem{theorem}[proposition]{Theorem}
\newtheorem{remark}[proposition]{Remark}

\title{Crank--Nicolson, the $\theta$-scheme,\\and Rannacher time-stepping}
\author{Roberto Cepeda Rocandio}
\date{May 2026; revised October 2026}

\begin{document}
\maketitle

\section{Setting and motivation}

FTCS and BTCS (previous two notes) are first-order in time: their local
truncation error is $O(\Delta\tau) + O(\Delta x^2)$. To balance the
spatial and temporal contributions one needs $\Delta\tau \sim
\Delta x^2$, which forces $M \sim N^2$ and total cost $O(N^3)$ for
both schemes (FTCS by CFL, BTCS by accuracy).

This note implements \emph{Crank--Nicolson} (CN), the unique scheme
in the convex family of FTCS and BTCS that is
\emph{second-order in time}: $O(\Delta\tau^2) + O(\Delta x^2)$. With
$M \sim N$, CN achieves $O(\Delta x^2)$ precision at $O(N^2)$ total
cost --- a factor-$N$ improvement over FTCS/BTCS in the same
precision regime. CN is the workhorse of finite-difference option
pricers in production.

The deliverable is a $\theta$-scheme infrastructure
(\texttt{quantlib.theta\_scheme}) that subsumes FTCS, BTCS and CN as
parameter choices, plus the CN pricer with Rannacher time-stepping
(\texttt{quantlib.cn}) to suppress spurious oscillations near the
strike kink.

% =====================================================================
\section{The $\theta$-family}
% =====================================================================

The three schemes FTCS, BTCS and CN are particular cases of a single
parametric family:
\begin{equation}
\frac{V_j^{n+1} - V_j^n}{\Delta\tau}
\;=\;
\theta\,\mathcal{L}_h V_j^{n+1} \,+\, (1-\theta)\,\mathcal{L}_h V_j^n,
\qquad \theta \in [0, 1],
\label{eq:theta-scheme}
\end{equation}
where $\mathcal{L}_h$ is the discrete spatial operator
\[
\mathcal{L}_h V_j
\;=\; \frac{\sigma^2}{2}\,\frac{V_{j+1} - 2V_j + V_{j-1}}{\Delta x^2}
   + \mu\,\frac{V_{j+1} - V_{j-1}}{2\Delta x}
   - r\,V_j.
\]
The cases of interest are
\begin{center}
\begin{tabular}{ccccc}
\hline
$\theta$ & Name & Type & Stability & Time order \\
\hline
$0$    & FTCS & Explicit & Conditional ($\alpha \le 1/2$) & $O(\Delta\tau)$ \\
$1/2$  & \textbf{Crank--Nicolson} & Implicit & Unconditional & $O(\Delta\tau^2)$ \\
$1$    & BTCS & Implicit & Unconditional & $O(\Delta\tau)$ \\
\hline
\end{tabular}
\end{center}
The order jump at $\theta = 1/2$ is the central observation of this
note. It is the pay-off of the symmetry between the new and old
levels: when the right-hand side is the average of $\mathcal{L}_h V$
at $\tau_n$ and $\tau_{n+1}$, the discrete equation is centred at the
\emph{midpoint} $\tau_{n+1/2}$, and the midpoint rule is one order
more accurate than the endpoint rules.

% =====================================================================
\section{Crank--Nicolson explicitly}
% =====================================================================

Setting $\theta = 1/2$ in \eqref{eq:theta-scheme}, expanding $\mathcal{L}_h$,
and rearranging $V^{n+1}$ to the left, $V^n$ to the right:
\begin{equation}
\boxed{\;
\beta_-\,V_{j-1}^{n+1} + \beta_0\,V_j^{n+1} + \beta_+\,V_{j+1}^{n+1}
\;=\;
\gamma_-\,V_{j-1}^n + \gamma_0\,V_j^n + \gamma_+\,V_{j+1}^n.
\;}
\label{eq:cn-stencil}
\end{equation}
With $\alpha = \frac{\sigma^2}{2}\,\Delta\tau/\Delta x^2$ and
$\nu = \mu\,\Delta\tau/(2\Delta x)$ as before, the LHS coefficients
are
\[
\beta_- = -\tfrac{1}{2}\alpha + \tfrac{1}{2}\nu, \quad
\beta_0 = 1 + \alpha + \tfrac{1}{2}r\Delta\tau, \quad
\beta_+ = -\tfrac{1}{2}\alpha - \tfrac{1}{2}\nu,
\]
and the RHS coefficients are
\[
\gamma_- = +\tfrac{1}{2}\alpha - \tfrac{1}{2}\nu, \quad
\gamma_0 = 1 - \alpha - \tfrac{1}{2}r\Delta\tau, \quad
\gamma_+ = +\tfrac{1}{2}\alpha + \tfrac{1}{2}\nu.
\]
Symmetry: in terms of the BTCS coefficients $b_\pm, b_0$ and the FTCS
coefficients $a_\pm, a_0$ of the previous notes, $\beta_\pm = \tfrac12 b_\pm$,
$\beta_0 - 1 = \tfrac12(b_0 - 1)$, and $\gamma_\pm = \tfrac12 a_\pm$,
$\gamma_0 - 1 = \tfrac12(a_0 - 1)$.
CN is, structurally, ``half BTCS plus half FTCS''.

% =====================================================================
\section{Why CN is second-order in time: the midpoint rule}
\label{sec:second-order}
% =====================================================================

The non-trivial claim is that the local truncation error of
\eqref{eq:cn-stencil} is $O(\Delta\tau^2) + O(\Delta x^2)$, despite
the fact that the time difference $(V^{n+1} - V^n)/\Delta\tau$ looks
like the same first-order forward Euler quotient that FTCS uses. The
gain comes from \emph{which time level the equation is centred on}.

Expand $V$ around the midpoint $\tau_{n+1/2} = \tau_n + \Delta\tau/2$:
\begin{align*}
V(\tau_{n+1}) &= V_* + \tfrac{1}{2}\Delta\tau\, V'_* + \tfrac{1}{8}\Delta\tau^2\, V''_*
+ \tfrac{1}{48}\Delta\tau^3\, V'''_* + O(\Delta\tau^4), \\
V(\tau_n)     &= V_* - \tfrac{1}{2}\Delta\tau\, V'_* + \tfrac{1}{8}\Delta\tau^2\, V''_*
- \tfrac{1}{48}\Delta\tau^3\, V'''_* + O(\Delta\tau^4),
\end{align*}
where $V_* := V(\tau_{n+1/2})$, $V'_* := V_\tau(\tau_{n+1/2})$, etc.
Subtracting and dividing by $\Delta\tau$:
\begin{equation}
\frac{V(\tau_{n+1}) - V(\tau_n)}{\Delta\tau}
\;=\; V_\tau(\tau_{n+1/2})
\;+\; \frac{\Delta\tau^2}{24}\,V_{\tau\tau\tau}(\tau_{n+1/2})
\;+\; O(\Delta\tau^4).
\label{eq:midpoint-time}
\end{equation}
The $\Delta\tau$ term \emph{cancels by symmetry}: the quotient is a
second-order approximation to $V_\tau$ at the midpoint.

Analogously, adding the two expansions of $\mathcal{L}_h V$ (one at
$\tau_n$, one at $\tau_{n+1}$) and dividing by 2:
\begin{equation}
\tfrac{1}{2}\!\left[\mathcal{L}_h V(\tau_{n+1}) + \mathcal{L}_h V(\tau_n)\right]
\;=\;
\mathcal{L}_h V(\tau_{n+1/2}) \,+\, O(\Delta\tau^2) \,+\, O(\Delta x^2),
\label{eq:midpoint-space}
\end{equation}
where the $O(\Delta\tau^2)$ comes from the time interpolation (linear
average of the endpoints) and the $O(\Delta x^2)$ is the spatial
truncation of $\mathcal{L}_h$. Combining \eqref{eq:midpoint-time} and
\eqref{eq:midpoint-space}, the CN equation
$(V^{n+1} - V^n)/\Delta\tau = \frac12(\mathcal{L}_h V^{n+1} + \mathcal{L}_h V^n)$
is the discretisation of the exact PDE $V_\tau = \mathcal{L} V$
\emph{at the midpoint}, with truncation
\[
\tau_h = O(\Delta\tau^2) + O(\Delta x^2).
\]
FTCS and BTCS, by contrast, evaluate the right-hand side at $\tau_n$ or
$\tau_{n+1}$ (not the midpoint); the time error there is $O(\Delta\tau)$.

% =====================================================================
\section{The linear system and matrix structure}
% =====================================================================

Equation \eqref{eq:cn-stencil} written for the interior nodes $j =
1, \dots, N-1$ becomes
\[
A\,\mathbf{V}^{n+1}_{\mathrm{int}}
\;=\; B\,\mathbf{V}^n_{\mathrm{int}}
   + \mathbf{b}_{\mathrm{bnd}}^n,
\]
where $A$ and $B$ are tridiagonal $(N-1)\times(N-1)$ matrices with
$(\beta_-, \beta_0, \beta_+)$ and $(\gamma_-, \gamma_0, \gamma_+)$ on
their diagonals, and the boundary vector mixes both time levels:
\begin{equation}
\mathbf{b}_{\mathrm{bnd}}^n
= \begin{pmatrix}
\gamma_-\,V_0^n - \beta_-\,V_0^{n+1} \\
0 \\
\vdots \\
0 \\
\gamma_+\,V_N^n - \beta_+\,V_N^{n+1}
\end{pmatrix}.
\label{eq:bnd-cn}
\end{equation}
Both lateral coefficients of the stencils contribute to the boundary
correction, in contrast to BTCS where only the LHS coefficients
appeared.

The per-step algorithm has three phases:
\begin{enumerate}[leftmargin=2em]
\item Apply the RHS stencil $B$ to $\mathbf{V}^n_{\mathrm{int}}$ as
a three-point stencil sweep ($O(N)$, no need to assemble $B$
explicitly).
\item Add the boundary correction $\mathbf{b}_{\mathrm{bnd}}^n$
($O(1)$: only two endpoints are non-zero).
\item Solve $A\,\mathbf{V}^{n+1}_{\mathrm{int}} = (\text{result})$ by
Thomas with the pre-factored $A$ ($O(N)$).
\end{enumerate}
Total cost per step: $O(N)$. Total cost: $O(N \cdot M)$, the same
as BTCS.

\paragraph{Pre-factoring $A$.} The matrix $A$ is constant in time
(again thanks to the constant-coefficient form of the transformed
PDE). Factoring $A$ once at setup and applying the factor at every
step is the standard pattern, identical to BTCS.

% =====================================================================
\section{Stability of the $\theta$-scheme}
% =====================================================================

The von Neumann amplification factor of \eqref{eq:theta-scheme},
ignoring the reactive term, is
\begin{equation}
g(\xi) = \frac{1 - 4\alpha(1-\theta)\sin^2(\xi/2) + i\,(1-\theta)\beta_C\sin\xi}
              {1 + 4\alpha\theta\sin^2(\xi/2) - i\,\theta\beta_C\sin\xi},
\label{eq:theta-amplification}
\end{equation}
with $\beta_C := \mu\Delta\tau/\Delta x$, $\xi \in [-\pi,\pi]$ the
Fourier wavenumber.

\begin{proposition}[Stability of the $\theta$-scheme]
The $\theta$-scheme is unconditionally von Neumann stable if and only
if $\theta \ge 1/2$. For $\theta < 1/2$, $|g| \le 1$ requires the
conditional bound
\[
\alpha \;\le\; \frac{1}{2(1 - 2\theta)}.
\]
In particular, $\theta = 0$ gives $\alpha \le 1/2$ (the FTCS CFL
condition), and $\theta = 1/2$ is the smallest $\theta$ for which the
bound is vacuous.
\end{proposition}

\begin{proof}[Sketch]
Set $s = \sin^2(\xi/2)$, $u = \sin\xi$, so $u^2 = 4s(1-s)$. Then
$|g|^2 = N(s,u)/D(s,u)$ where both numerator and denominator are
quadratic in $\alpha$. The condition $|g|^2 \le 1$ becomes $D - N \ge
0$, which after expansion is
\[
8\alpha s\,\bigl[1 + 2\alpha(2\theta - 1)s\bigr]
\,+\,(2\theta - 1)\,\beta_C^2 u^2
\,\ge\, 0.
\]
For $\theta \ge 1/2$, the bracket is at least $1$ and the $\beta_C$ term
is non-negative, for any $\alpha \ge 0$, $s \in [0,1]$. For
$\theta < 1/2$, take the Nyquist mode $s = 1$, $u = 0$: the condition
becomes $1 - 2\alpha(1 - 2\theta) \ge 0$, which is the stated
inequality. (Full algebra in Strikwerda, Chapter~6, or \cite{Smith1985}.) \qed
\end{proof}

CN at $\theta = 1/2$ is therefore \emph{exactly on the boundary}
between conditional and unconditional stability. This is the structural
reason why CN, although stable, has weaker damping properties than
BTCS, as we now examine.

% =====================================================================
\section{The kink-induced oscillation problem}
\label{sec:kink}
% =====================================================================

The amplification factor at the Nyquist wavenumber $\xi = \pi$
(highest representable frequency on the grid, $\sin^2(\pi/2) = 1$,
$\sin\pi = 0$) is
\[
g_\theta(\pi) = \frac{1 - 4\alpha(1-\theta)}{1 + 4\alpha\theta}.
\]
For BTCS ($\theta = 1$) and CN ($\theta = 1/2$):
\[
g_{\mathrm{BTCS}}(\pi) = \frac{1}{1 + 4\alpha}
\xrightarrow{\alpha \to \infty} 0,
\qquad
g_{\mathrm{CN}}(\pi) = \frac{1 - 2\alpha}{1 + 2\alpha}
\xrightarrow{\alpha \to \infty} -1.
\]
\textbf{BTCS damps the Nyquist mode exponentially; CN preserves it
with a sign flip.} In the regime of interest for CN ($\Delta\tau \sim
\Delta x$, so $\alpha \sim 1/\Delta x \to \infty$ as $N \to \infty$),
$|g_{\mathrm{CN}}(\pi)|$ is essentially $1$, and the Nyquist mode is
preserved indefinitely while alternating sign in time.

This is harmless for smooth initial data: the Fourier coefficient of
the Nyquist mode in the initial vector $V^0$ is exponentially small,
and ``preserving exponentially small'' yields nothing visible.

\paragraph{The payoff is not smooth.} A vanilla call payoff
$\phi(x) = K(e^x - 1)^+$ has a kink at $x = 0$: $\phi$ is continuous
but $\phi'$ has a jump. The Fourier expansion of a function with such
a kink has coefficients decaying only as $1/k^2$, not exponentially.
The Nyquist coefficient in $V^0$ is therefore $O(1/N^2)$ --- not
exponentially small.

\paragraph{Consequence.} An $O(1/N^2)$ amplitude at the Nyquist mode,
preserved across $M$ time steps with sign flip $g \approx -1$ at each,
produces a visible oscillation pattern in the computed $V^M$:
\begin{itemize}[leftmargin=2em]
\item Bounded ($|g| \le 1$, so the oscillation does not grow).
\item Persistent (each step preserves the magnitude).
\item Localised near the kink (where the kink-induced mode lives).
\item Alternating sign at consecutive grid nodes (the spatial
signature of $\xi = \pi$).
\end{itemize}
The visible artefact is a sawtooth of small amplitude near the
strike, distinct from the explosive sawtooth of FTCS instability
(see the FTCS note) by being bounded but visible. How much it
degrades the price depends on $\Delta\tau/\Delta x$. With $\Delta\tau$
comparable to $\Delta x$ the price at the strike is barely affected on
the grids used here: for the at-the-money call of \texttt{validate\_cn.py}
($S = K = 100$, $\sigma = 0.2$, $T = 1$) with $M = N/2$
($\Delta\tau/\Delta x = 1.25$), vanilla CN converges with error ratio
$4.00$ per refinement ($5.68\times10^{-3}$ at $N = 100$ to
$8.86\times10^{-5}$ at $N = 800$), slightly more accurate than with
Rannacher. The loss of order appears when $\Delta\tau \gg \Delta x$:
test~[3] of \texttt{validate\_cn.py} ($N = 800$, $M = 10, 20, 40$) gives
error ratios $2.11$ and $2.53$ under halving of $\Delta\tau$, against
$3.89$ and $3.55$ with Rannacher. Giles and Carter~\cite{GilesCarter2006}
observe first-order convergence of the price at the largest
$\Delta\tau/\Delta x$ they test, and note that for smaller ratios the
same behaviour shows only on much finer grids. The Greeks (gamma in
particular) inherit the oscillation.

% =====================================================================
\section{Rannacher time-stepping}
% =====================================================================

The standard fix goes back to Rannacher~\cite{Rannacher1984}, who
replaced the first CN time step by two half-time-step implicit (BTCS)
steps. The version used here, from Giles and
Carter~\cite{GilesCarter2006}, replaces the first \emph{two} CN time
steps by \emph{four} half-time-step BTCS steps, each of length
$\Delta\tau/2$:
\begin{equation}
\underbrace{\mathrm{BTCS}_{\Delta\tau/2}\circ\mathrm{BTCS}_{\Delta\tau/2}\circ\mathrm{BTCS}_{\Delta\tau/2}\circ\mathrm{BTCS}_{\Delta\tau/2}}_{\text{4 half-steps, total time }2\Delta\tau}
\;\circ\;
\underbrace{\mathrm{CN}_{\Delta\tau} \circ \cdots \circ \mathrm{CN}_{\Delta\tau}}_{M-2\text{ full steps}}.
\label{eq:rannacher}
\end{equation}
\paragraph{Why it works.} Each BTCS half-step damps the Nyquist mode
by a factor $g_{\mathrm{BTCS},\Delta\tau/2}(\pi) = 1/(1 + 2\alpha)$
($\alpha$ is the CN Fourier number; the half-step uses $\alpha/2$,
giving $1/(1 + 2\alpha)$). After four half-steps the Nyquist mode is
multiplied by $1/(1+2\alpha)^4$, which for $\alpha = O(N) \gg 1$ is
$O(1/N^4)$ --- well below the $O(\Delta x^2) = O(1/N^2)$ floor of the
discretisation. By the time CN takes over, the kink-induced
high-frequency modes are damped to round-off, and CN proceeds with
effectively smooth data.

\paragraph{Why two full steps.} Rannacher's original start-up (one CN
step replaced by two implicit half-steps) already recovers
second-order convergence in $L^2$. Giles \&
Carter~\cite{GilesCarter2006} show that replacing two CN steps by four
half-steps is needed for second-order accuracy in $L^\infty$ for the
value, delta and gamma; that is the protocol \eqref{eq:rannacher} and
the default here.

\paragraph{Why BTCS, not finer CN.} Refining $\Delta\tau$ does remove
the artefact, but only once $\alpha = \sigma^2\Delta\tau/(2\Delta x^2)$
is $O(1)$. As $\Delta\tau \to 0$ at fixed $\Delta x$, $\alpha \to 0$
and $g_{\mathrm{CN}}(\pi) = (1 - 2\alpha)/(1 + 2\alpha) \to +1$ with no
sign flip: over $[0, T]$ the Nyquist mode decays like
$e^{-2\sigma^2 T/\Delta x^2}$, as it should. For $\alpha \le 1/2$ there
is no sign flip at all, and for $\alpha$ of order one the flipping mode
loses a factor $(2\alpha - 1)/(2\alpha + 1)$ per step ($1/3$ at
$\alpha = 1$). The persistent sign-flipping mode
needs $\alpha \gg 1$, i.e.\ $\Delta\tau$ large relative to $\Delta x^2$,
which is exactly the regime $\Delta\tau \sim \Delta x$ in which CN is
worth using. Taking $\Delta\tau \sim \Delta x^2$ instead would bring back
$M \sim N^2$ and the $O(N^3)$ cost of FTCS/BTCS. Within the CN regime
only a dissipative scheme can remove the irregular high-frequency
content, and BTCS is the natural choice because the same Thomas
factorisation can be reused with the appropriate coefficients.

\paragraph{Cost.} The four BTCS half-steps replace two CN steps, so the
protocol adds two tridiagonal solves, about $2/M$ of the CN cost, plus
one extra factorisation for the half-step matrix: a noticeable overhead
at $M = 10$, negligible by $M = 100$, and small compared to the
convergence gain when $\Delta\tau \gg \Delta x$.

% =====================================================================
\section{Implementation outline}
% =====================================================================

The implementation consists of two modules:
\begin{itemize}[leftmargin=2em]
\item \texttt{quantlib.theta\_scheme} (Python) /
\texttt{quant::pde::theta\_scheme\_*} (C++): a generic $\theta$-stepper
parametrised by $\theta$. Internally pre-factors the LHS matrix and
runs a configurable number of time steps. Reusable as scaffolding
for the American put (where the LHS solve is replaced by PSOR).
\item \texttt{quantlib.cn} (Python) /
\texttt{quant::pde::cn\_*} (C++): the user-facing CN pricer for
European options, with Rannacher smoothing enabled by default
(\texttt{rannacher\_steps=2}). Setting \texttt{rannacher\_steps=0}
gives vanilla CN, exposed to make the oscillation phenomenon
testable.
\end{itemize}

The signature mirrors BTCS:
\begin{verbatim}
cn_european_call(S, K, r, sigma, T, *, N, M, n_sigma=4.0,
                 rannacher_steps=2)
\end{verbatim}

The $\theta$-scheme module exposes a low-level kernel
\texttt{theta\_march(grid, V0, theta, bc\_lower, bc\_upper)} that the
CN pricer calls twice for the Rannacher protocol: once with
$\theta = 1$ for the initial half-stepped BTCS warm-up, then with
$\theta = 1/2$ for the rest.

\paragraph{No replacement of FTCS/BTCS.} The FTCS and BTCS modules
do not use it. The $\theta$-scheme infrastructure is a
parallel mechanism intended for CN and any future generalisations.

% =====================================================================
\section{Validation strategy}
% =====================================================================

Six tests in \texttt{validate\_cn.py} certify the new pricer
(\texttt{test\_cn.cpp} repeats the others for the C++ pricer; the fourth
is in \texttt{test\_theta\_scheme.cpp}):
\begin{enumerate}[leftmargin=2em]
\item Cross-validation against Black--Scholes closed form on a
parameter grid (the standard battery used for FTCS and BTCS): calls
and puts on eight parameter sets at $N = 400$, $M = 200$, tolerance
$5\times10^{-3}$.
\item \textbf{Quadratic convergence in time:} at $N = 800$, halving
$\Delta\tau$ ($M = 10, 20, 40$) must reduce the error by a factor
above $3$, against $\approx 2$ for BTCS; the measured ratios are $3.89$
and $3.55$. This is \emph{the} test of CN.
\item \textbf{Vanilla CN (without Rannacher) has degraded convergence:}
the same problem with \texttt{rannacher\_steps=0} must give ratios
below $3$ (measured $2.11$ and $2.53$) and an error at $M = 10$ at least
five times the Rannacher one. This is the empirical evidence of the
kink-induced oscillation.
\item Reproduction of BTCS via $\theta = 1$ in the $\theta$-scheme:
\texttt{theta\_march} with $\theta = 1$ matches the BTCS pricer to
round-off.
\item Put--call parity at five spots.
\item Input validation ($S \le 0$, negative \texttt{rannacher\_steps},
\texttt{rannacher\_steps} $> M$, spot outside the truncated domain).
\end{enumerate}

The third test is the operational signature of the kink-induced
oscillation, analogous to the CFL-violation diagnostic of FTCS: a
configuration that produces visibly degraded output, captured as a
passing assertion to prevent regression.

% =====================================================================
\section*{References}
% =====================================================================

\begin{thebibliography}{9}

\bibitem{TavellaRandall2000}
D. Tavella and C. Randall.
\emph{Pricing Financial Instruments: The Finite Difference Method.}
Wiley, 2000. Chapter~2 (Analysis of Finite Difference Methods).

\bibitem{Seydel2017}
R. Seydel.
\emph{Tools for Computational Finance.}
6th ed., Springer, 2017. Chapter~4 (Standard Methods for Standard
Options): finite-difference schemes, including CN, and their accuracy.

\bibitem{Strikwerda2004}
J.~C. Strikwerda.
\emph{Finite Difference Schemes and Partial Differential Equations.}
2nd ed., SIAM, 2004. \S6.3 (stability analysis of the $\theta$-method).

\bibitem{Smith1985}
G.~D. Smith.
\emph{Numerical Solution of Partial Differential Equations:
Finite Difference Methods.}
3rd ed., OUP, 1985. The $\theta$-family of schemes.

\bibitem{Rannacher1984}
R. Rannacher.
\emph{Finite element solution of diffusion problems with irregular data.}
Numerische Mathematik 43, 309--327, 1984. The original paper
introducing the smoothing protocol (one CN step replaced by two
implicit half-steps).

\bibitem{GilesCarter2006}
M.~B. Giles and R. Carter.
\emph{Convergence analysis of Crank-Nicolson and Rannacher
time-marching.}
Journal of Computational Finance 9(4), 89--112, 2006. Results for
option pricing: Rannacher's original start-up gives second order in
$L^2$; replacing two CN steps by four implicit half-steps gives
second-order $L^\infty$ accuracy for the value, delta and gamma.

\end{thebibliography}

\end{document}
